\chapter{Logic and quantifiers}\label{ch:logic}
\index{logic}

Every claim in this book is either a statement you will prove or a statement you will
refute. Before any technique, you need the grammar those statements are written in:
connectives that combine simple claims, and quantifiers that say \emph{how many} objects
a claim covers. After this chapter you will be able to write any contest-level claim in
symbols, negate it mechanically without thinking, and say precisely what a proof or a
counterexample of it would have to deliver. Nothing here is deep; everything here gets
used weekly. Read it once carefully and it will save you from the most common proof
errors in the rest of the book.

% from-zero:why:begin
\paragraph{Why this exists.} Computers do exactly what they are told, so every claim about a program
(``this sort always finishes'') needs a justification stronger than ``I tried it''. Logic is the language
in which such claims are stated unambiguously, and proof is how they are settled. Olympiad problems ask for
the same thing: not an answer, but an argument that holds in every case. Learn the language first; the proof
techniques of the next chapters are written in it.
% from-zero:why:end

% from-zero:begin
\begin{rulebox}[Before you start]
Nothing. If you can add, multiply and tell an even number from an odd one, you can start. The \emph{Start here} chapter lists the few school-level facts we use (what an integer is, what ``divides'' means, what a prime is) and takes two minutes to skim.
\end{rulebox}
\begin{wins}[After this chapter you can]
\begin{itemize}
\item decide whether a sentence is a statement with a definite truth value;
\item build a truth table and use it to check that two statements say the same thing;
\item write ``every'' and ``some'' claims with quantifiers and negate them correctly;
\item say exactly what a proof or a counterexample of a claim has to show.
\end{itemize}
\end{wins}
\begin{principle}[How to think about this]
Mathematical language is just English with the ambiguity removed. When you read a symbolic statement, say it aloud in plain words; when you meet an English claim, ask ``for \emph{every} what? for \emph{some} what?'' Almost every proof mistake starts with a quantifier that was quietly misread.
\end{principle}
% from-zero:end

\section{Statements and connectives}\label{sec:statements}

\begin{definition}[Statement]\index{statement}\index{proposition|see{statement}}
A \emph{statement} (also called a \emph{proposition}) is a sentence that is either
true or false. It need not be provable, known, or interesting --- it only needs a
definite truth value.
\end{definition}

``7 is prime'' is a statement, and a true one. ``9 is prime'' is a statement, and a
false one. ``Every even integer greater than $2$ is a sum of two primes'' is a
statement --- nobody knows which of the two truth values it has, but it has one.
On the other hand ``$x+2=7$'' is not a statement, because it is true for some $x$
and false for others; it is an \emph{open formula}\index{open formula}, and
Section~\ref{sec:quantifiers} shows how quantifiers close it. Questions and
commands (``is it raining?'', ``solve this'') are not statements at all.

Statements combine into larger ones using \emph{connectives}\index{connective}.

\begin{definition}[Basic connectives]\index{conjunction}\index{disjunction}\index{negation}
Let $P$ and $Q$ be statements.
\begin{itemize}
\item $P \land Q$ (\emph{conjunction}, ``$P$ and $Q$'') is true exactly when both are true.
\item $P \lor Q$ (\emph{disjunction}, ``$P$ or $Q$'') is true exactly when at least one
is true --- \textbf{including when both are true}.
\item $\lnot P$ (\emph{negation}, ``not $P$'') is true exactly when $P$ is false.
\end{itemize}
\end{definition}

\begin{center}
\begin{tabular}{cc|ccc}
\toprule
$P$ & $Q$ & $P\land Q$ & $P\lor Q$ & $\lnot P$\\
\midrule
T & T & T & T & F\\
T & F & F & T & F\\
F & T & F & T & T\\
F & F & F & F & T\\
\bottomrule
\end{tabular}
\end{center}

\begin{remark}\index{or!inclusive}
Mathematics uses the \emph{inclusive} or: ``$P$ or $Q$'' allows both. English does
not always. ``You may take the exam in June or in August'' usually means one or the
other. ``A student who studied or who is lucky will pass'' allows both. When a
contest statement wants the exclusive version it will say so: ``exactly one of''.
\end{remark}

\section{Truth tables}\label{sec:truthtables}
\index{truth table}

A connective is defined by what it does to truth values, so the whole behaviour of a
compound statement is computable line by line. Writing out all $2^n$ rows for $n$
atomic statements is called a \emph{truth table}. It is the brute-force verifier of
propositional logic: two formulas are \emph{equivalent}\index{equivalence!of formulas}
--- interchangeable inside any larger statement --- exactly when their truth tables
agree on every row.

Here is the most useful equivalence there is, in action:

\begin{proposition}[De Morgan's laws for statements]\label{prop:demorgan}
\index{De Morgan's laws}
$\lnot(P\land Q)$ is equivalent to $\lnot P \lor \lnot Q$, and
$\lnot(P\lor Q)$ is equivalent to $\lnot P \land \lnot Q$.
\end{proposition}

\begin{proof}
Compare the two sides of the first law on all four cases:
\begin{center}
\begin{tabular}{cc|cc|c|cc}
\toprule
$P$ & $Q$ & $P\land Q$ & $\lnot(P\land Q)$ & $\lnot P\lor\lnot Q$ & $\lnot P$ & $\lnot Q$\\
\midrule
T&T&T&F&F&F&F\\
T&F&F&T&T&F&T\\
F&T&F&T&T&T&F\\
F&F&F&T&T&T&T\\
\bottomrule
\end{tabular}
\end{center}
Columns four and five agree on every row, so the formulas are equivalent. The second
law is checked the same way.
\end{proof}

Read De Morgan in English: ``not both'' means ``at least one fails''; ``neither''
means ``both fail''. You will use it every time you negate a compound statement in
Section~\ref{sec:negation}.

\begin{exercise}[easy]\label{exr:tt-notor}
Write the truth table of $\lnot P \lor Q$ for all four rows of $(P,Q)$. Keep it ---
in Section~\ref{sec:implication} this formula gets a name.
\end{exercise}

\section{Implication and its family}\label{sec:implication}
\index{implication}

Most theorems have the shape ``if $P$ then $Q$''.

\begin{definition}[Implication]\index{hypothesis}\index{conclusion}
$P \Rightarrow Q$ (\emph{implication}) is false in exactly one case: $P$ true and $Q$
false. $P$ is the \emph{hypothesis}, $Q$ the \emph{conclusion}. English renders it as
``if $P$ then $Q$'', ``$P$ implies $Q$'', ``$Q$ if $P$'', ``$P$ only if $Q$'', and
``$Q$ whenever $P$''.
\end{definition}

The last two rows of the table are the part everyone trips on:

\begin{center}
\begin{tabular}{cc|c}
\toprule
$P$ & $Q$ & $P\Rightarrow Q$\\
\midrule
T&T&T\\
T&F&F\\
F&T&T\\
F&F&T\\
\bottomrule
\end{tabular}
\end{center}

\begin{remark}\index{vacuous truth}
If $P$ is false, $P \Rightarrow Q$ is automatically true --- \emph{vacuously} true.
``If $2+2=5$ then this textbook has a million pages'' is a true statement, and proves
nothing about the book. An implication only ever makes a promise about the case where
the hypothesis holds.
\end{remark}

Exercise~\ref{exr:tt-notor} now pays off: the table you computed shows that
$P \Rightarrow Q$ is equivalent to $\lnot P \lor Q$. This one fact explains the whole
implication family and, later, how to negate an implication.

\begin{definition}[The implication family]\index{converse}\index{inverse}\index{contrapositive}
For $P \Rightarrow Q$:
\begin{itemize}
\item the \emph{converse} is $Q \Rightarrow P$;
\item the \emph{inverse} is $\lnot P \Rightarrow \lnot Q$;
\item the \emph{contrapositive} is $\lnot Q \Rightarrow \lnot P$.
\end{itemize}
\end{definition}

\begin{figure}[ht]\centering
% Figure: the implication family. Diagonal pairs are equivalent;
% horizontal/vertical neighbours are different claims.
% Source for Figure~\ref{fig:family}. TikZ only; regenerable by compiling main.tex.
\begin{tikzpicture}[
  box/.style={draw,inner sep=5pt,font=\small,rounded corners=1pt},
  equiv/.style={double,double distance=1.5pt,-,accent},
  diff/.style={-,dashed,rulegray},
  lab/.style={font=\footnotesize\sffamily,text=slate},
]
\node[box] (imp)  at (0,2.2) {$P \Rightarrow Q$};
\node[box] (conv) at (5.4,2.2) {$Q \Rightarrow P$};
\node[box] (inv)  at (0,0) {$\lnot P \Rightarrow \lnot Q$};
\node[box] (contra) at (5.4,0) {$\lnot Q \Rightarrow \lnot P$};
\node[lab,anchor=south] at (0,2.45) {implication};
\node[lab,anchor=south] at (5.4,2.45) {converse};
\node[lab,anchor=north] at (0,-0.25) {inverse};
\node[lab,anchor=north] at (5.4,-0.25) {contrapositive};
\draw[equiv] (imp) -- (contra) node[midway,lab,sloped,above=-0.5pt]{same truth value};
\draw[equiv] (conv) -- (inv) node[midway,lab,sloped,above=-0.5pt]{same truth value};
\draw[diff] (imp) -- (conv) node[midway,lab,above]{different claims};
\draw[diff] (imp) -- (inv) node[midway,lab,left]{different claims};
\end{tikzpicture}

\caption{The implication family. Notice that equivalence lives on the
\emph{diagonals}: an implication and its contrapositive are the same claim
rephrased, and so are the converse and the inverse. Neighbours are different claims.}
\label{fig:family}
\end{figure}

\begin{theorem}\label{thm:contrapositive}
$P \Rightarrow Q$ is equivalent to its contrapositive $\lnot Q \Rightarrow \lnot P$,
and to $\lnot P \lor Q$. It is \emph{not} equivalent to its converse
$Q \Rightarrow P$.
\end{theorem}

\begin{proof}
Apply the equivalence $X \Rightarrow Y \equiv \lnot X \lor Y$ twice, with
commutativity of $\lor$ and double negation between them:
$P \Rightarrow Q \equiv \lnot P \lor Q \equiv Q \lor \lnot P
\equiv \lnot(\lnot Q) \lor \lnot P \equiv \lnot Q \Rightarrow \lnot P$.
For the converse, the rows $(P,Q)=(T,F)$ and $(F,T)$ of the truth table already
disagree between $P\Rightarrow Q$ and $Q\Rightarrow P$, so no equivalence is possible.
\end{proof}

The contrapositive is the single most useful free move in proof-writing: whenever a
statement is awkward in the form given, its contrapositive is the same statement in a
different shape. Chapter~2 exploits this constantly.

\begin{definition}[Biconditional]\index{biconditional}\index{iff}
$P \iff Q$ (``$P$ if and only if $Q$'', written \emph{iff}) is
$(P \Rightarrow Q) \land (Q \Rightarrow P)$: true when $P$ and $Q$ have the same
truth value. A claim stated as iff is \emph{two} obligations, not one.
\end{definition}

\begin{example}\label{ex:ttree}
Evaluate $\lnot (P \land Q) \Rightarrow R$ when $P$ and $Q$ are true and $R$ is false.
\par\smallskip
\textbf{Strategy.} Evaluate inside out: leaves first, then each connective, never
reasoning about the whole formula at once. A tree of subformulas tracks the order.
\par\smallskip
\begin{solution}
$P \land Q$ is T since both inputs are T. Then $\lnot(P\land Q)$ is F, so the
statement is $F \Rightarrow F$, which is T. Figure~\ref{fig:tt-tree} shows the same
computation for the companion formula $(P\land Q) \Rightarrow R$, where this
assignment is the \emph{only} row that fails.
\end{solution}
\end{example}

\begin{figure}[ht]\centering
% Figure: evaluation tree for (P /\ Q) => R at the row P=T, Q=T, R=F.
% Source for Figure~\ref{fig:tt-tree}. TikZ only; regenerable by compiling main.tex.
\begin{tikzpicture}[
  level distance=10mm,
  level 1/.style={sibling distance=52mm},
  level 2/.style={sibling distance=26mm},
  every node/.style={font=\small},
  op/.style={draw,inner sep=2.5pt,rounded corners=1pt},
  leaf/.style={font=\small},
  val/.style={font=\footnotesize\sffamily,text=slate},
  badval/.style={font=\footnotesize\sffamily,text=amber},
  edge from parent/.style={draw=rulegray,-},
]
\node[op] (root) {$P \land Q \Rightarrow R$}
  child { node[op] (and) {$P \land Q$}
    child { node[leaf] (p) {$P$} }
    child { node[leaf] (q) {$Q$} }
  }
  child { node[leaf] (r) {$R$} };
% truth values propagated upward for the assignment P=T,Q=T,R=F
\node[val,anchor=west] at (p.east) {= T};
\node[val,anchor=west] at (q.east) {= T};
\node[badval,anchor=west] at (r.east) {= F};
\node[val,anchor=east] at (and.west) {= T};
\node[badval,anchor=south] at (root.north) {= F};
\node[draw=none,font=\footnotesize\itshape,text=slate] at (0,-2.9)
  {assignment $P{=}T,\ Q{=}T,\ R{=}F$: the only row where the implication fails};
\end{tikzpicture}

\caption{Evaluation tree of $(P\land Q)\Rightarrow R$ on the assignment
$P{=}T,\,Q{=}T,\,R{=}F$. Notice that the implication is false exactly where the left
child evaluates to T and the right to F --- one row out of eight decides it.}
\label{fig:tt-tree}
\end{figure}

\begin{exercise}[medium]\label{exr:converse}
For the statement ``if $n$ is divisible by $6$ then $n$ is even'', write the
converse, the inverse and the contrapositive, and decide which of the four
statements are true over the integers.
\end{exercise}

\section{Quantifiers}\label{sec:quantifiers}
\index{quantifier}

``$x+2=7$'' was not a statement because $x$ was floating free. Quantifiers bind such
variables.

\begin{definition}[Universal and existential quantifiers]\index{universal quantifier}%
\index{existential quantifier}\index{universe of discourse}
Fix a \emph{universe of discourse} $U$ --- the set of objects $x$ may range over. For
an open formula $P(x)$:
\begin{itemize}
\item $\forall x\, P(x)$ (\emph{for all}): true when $P(x)$ holds for every $x\in U$.
\item $\exists x\, P(x)$ (\emph{there exists}): true when $P(x)$ holds for at least
one $x\in U$.
\item $\exists! x\, P(x)$ (\emph{there exists a unique}): exactly one $x\in U$
satisfies $P(x)$.
\end{itemize}
Write $\forall x\in U$ or $\exists x\in U$ when the universe needs naming.
\end{definition}

Read the two quantifiers as jobs. $\forall$ says \emph{be ready for anything}:
to verify $\forall x\,P(x)$ you must handle every single $x$, and to refute it one
\emph{counterexample}\index{counterexample} suffices. $\exists$ says \emph{produce a
witness}\index{witness}: to verify it, hand over one $x$ that works; to refute it you
must exclude \emph{every} candidate. Asymmetry all the way down.

\begin{example}[$\forall\exists$ versus $\exists\forall$]\label{ex:quant-swap}
On the set of contestants $\{\text{Ana, Budi, Citra, Dewi}\}$ and problems
$\{P1,P2,P3,P4\}$, let $S(c,p)$ mean ``contestant $c$ solved problem $p$'', with
solved problems:
\begin{center}
\begin{tabular}{l|cccc}
\toprule
 & $P1$ & $P2$ & $P3$ & $P4$\\
\midrule
Ana & \checkmark & & &\\
Budi & & \checkmark & &\\
Citra & & & \checkmark &\\
Dewi & & & & \checkmark\\
\bottomrule
\end{tabular}
\end{center}
Compare $\forall c\,\exists p\, S(c,p)$ and $\exists p\,\forall c\, S(c,p)$.
\par\smallskip
\textbf{Strategy.} Read the first quantifier as the \emph{job}, the second as the
\emph{answer}. In $\forall c\,\exists p$, the answer $p$ may depend on $c$ --- each
contestant may bring a different solved problem. In $\exists p\,\forall c$, one
single $p$ must serve everyone.
\par\smallskip
\begin{solution}
Every row of the table has a checkmark, so $\forall c\,\exists p\,S(c,p)$ is
\emph{true}. No column is full of checkmarks, so $\exists p\,\forall c\,S(c,p)$ is
\emph{false}. The two claims differ: order of quantifiers is not decoration.
\end{solution}
\end{example}

\begin{figure}[ht]\centering
% Figure: forall-exists vs exists-forall on two 4x4 grids of marks.
% Source for Figure~\ref{fig:quant-order}. TikZ only; regenerable by compiling main.tex.
\begin{tikzpicture}[cell/.style={minimum size=5.6mm},mk/.style={circle,fill=accent,inner sep=1.6pt},
  lab/.style={font=\footnotesize\sffamily,text=slate}]
% left grid: marks on the diagonal -- every row has a mark, no column covers all rows
\foreach \r in {1,...,4} \foreach \c in {1,...,4}
  \node[cell,draw=rulegray] at (0.6*\c,-0.6*\r) {};
\foreach \i in {1,...,4} \node[mk] at (0.6*\i,-0.6*\i) {};
\foreach \r in {1,...,4} \node[lab,anchor=east] at (0.24,-0.6*\r) {$x_\r$};
\foreach \c in {1,...,4} \node[lab,anchor=south] at (0.6*\c,-0.26) {$y_\c$};
\node[lab] at (1.5,-2.85) {$\forall x\,\exists y$: T\qquad $\exists y\,\forall x$: F};

% right grid: one whole column marked -- a single column covers every row
\begin{scope}[xshift=58mm]
\foreach \r in {1,...,4} \foreach \c in {1,...,4}
  \node[cell,draw=rulegray] at (0.6*\c,-0.6*\r) {};
\foreach \r in {1,...,4} \node[mk] at (0.6*2,-0.6*\r) {};
\foreach \r in {1,...,4} \node[lab,anchor=east] at (0.24,-0.6*\r) {$x_\r$};
\foreach \c in {1,...,4} \node[lab,anchor=south] at (0.6*\c,-0.26) {$y_\c$};
\node[lab] at (1.5,-2.85) {$\forall x\,\exists y$: T\qquad $\exists y\,\forall x$: T};
\end{scope}
\end{tikzpicture}

\caption{Two finite worlds. Left: marks on the diagonal --- every row has a mark
($\forall x\,\exists y$ true) but no column covers all rows
($\exists y\,\forall x$ false). Right: one column fully marked, so both claims
hold. Notice that $\exists y\,\forall x$ is the \emph{stronger} claim: when it
holds, $\forall x\,\exists y$ follows.}
\label{fig:quant-order}
\end{figure}

The same swap runs under the hood of definitions you will meet later.
$\forall \varepsilon>0\,\exists \delta>0 \dots$ in analysis means ``however small the
tolerance, there is a response \emph{that may depend on it}''. Confusing the two
orders is the most common way definitions get misread; Listing~\ref{lst:verifyq}
lets you check both verdicts by brute force on any small world.

\begin{checkpoint}[Predict before you run]
In Listing~\ref{lst:verifyq}, world 1 is the table of
Example~\ref{ex:quant-swap}. Before running it, predict the four printed verdicts:
two for world 1, then world 2 (everyone solved $P1$ only) and world 3 (Dewi solved
nothing).
\end{checkpoint}

\begin{lstlisting}[language=Python,caption={\texttt{verify\_quantifiers.py}: checking
$\forall\exists$ vs $\exists\forall$ by brute force on finite worlds.},
label={lst:verifyq}]
def verdicts(passed):
    fe = all(any(p in passed[c] for p in problems)
             for c in contestants)
    ef = any(all(p in passed[c] for c in contestants)
             for p in problems)
    return fe, ef
\end{lstlisting}
The full script (in \texttt{code/verify\_quantifiers.py}) builds three worlds and
prints both verdicts for each:
\begin{lstlisting}[numbers=none]
world 1:
  every contestant solved some problem : True
  some problem was solved by everyone  : False
world 2:
  every contestant solved some problem : True
  some problem was solved by everyone  : True
world 3:
  every contestant solved some problem : False
  some problem was solved by everyone  : False
\end{lstlisting}
Notice that the first claim can hold without the second (world 1), both can hold
(world 2), and both can fail (world 3) --- the claims are strictly different, and
$\exists p\,\forall c$ implies $\forall c\,\exists p$ but not conversely.

\begin{example}[A finite world]\label{ex:tarski}
In the $4\times4$ grid of Figure~\ref{fig:tarski}, each cell holds one symbol from
$\{\triangle,\circ,\square\}$. Evaluate: (a) every row contains a $\triangle$;
(b) some column consists entirely of $\circ$; (c) every $\circ$ shares its row with
a $\square$; (d) some column contains both a $\triangle$ and a $\square$.
\par\smallskip
\textbf{Strategy.} When the universe is small and finite, work on the model, not on
the symbols. $\forall$ becomes a checklist over rows or cells; $\exists$ becomes a
search that stops at the first hit.
\par\smallskip
\begin{solution}
(a) Rows $1$--$4$ have $\triangle$ in columns $1,2,3,4$ respectively: true.
(b) Scanning columns $1$--$4$, each contains a non-$\circ$ cell: false.
(c) Every row contains a $\square$ (columns $3,4,1,2$), and every row contains
$\circ$ cells, so each $\circ$ shares its row with a $\square$: true.
(d) Column $2$ contains $\triangle$ (row $2$) and $\square$ (row $4$): true.
\end{solution}
\end{example}

\begin{figure}[ht]\centering
% Figure: the 4x4 finite world used in Worked Example 1.3.
% Source for Figure~\ref{fig:tarski}. TikZ only; regenerable by compiling main.tex.
\begin{tikzpicture}[cell/.style={minimum size=7mm},lab/.style={font=\footnotesize\sffamily,text=slate}]
\foreach \r in {1,...,4} \foreach \c in {1,...,4}
  \node[cell,draw=rulegray] at (0.75*\c,-0.75*\r) {};
% symbols, row by row:  triangle / circle / square
\node at (0.75*1,-0.75*1) {$\triangle$}; \node at (0.75*2,-0.75*1) {$\circ$};
\node at (0.75*3,-0.75*1) {$\square$};  \node at (0.75*4,-0.75*1) {$\circ$};
\node at (0.75*1,-0.75*2) {$\circ$};    \node at (0.75*2,-0.75*2) {$\triangle$};
\node at (0.75*3,-0.75*2) {$\circ$};    \node at (0.75*4,-0.75*2) {$\square$};
\node at (0.75*1,-0.75*3) {$\square$};  \node at (0.75*2,-0.75*3) {$\circ$};
\node at (0.75*3,-0.75*3) {$\triangle$};\node at (0.75*4,-0.75*3) {$\circ$};
\node at (0.75*1,-0.75*4) {$\circ$};    \node at (0.75*2,-0.75*4) {$\square$};
\node at (0.75*3,-0.75*4) {$\circ$};    \node at (0.75*4,-0.75*4) {$\triangle$};
\foreach \r in {1,...,4} \node[lab,anchor=east] at (0.30,-0.75*\r) {row \r};
\foreach \c in {1,...,4} \node[lab,anchor=south] at (0.75*\c,-0.36) {col \c};
\end{tikzpicture}

\caption{The finite world of Example~\ref{ex:tarski}. Notice that the $\triangle$
marks form a diagonal --- that is why (a) is true while every column fails the
``all $\circ$'' test in (b). Small worlds in this style go back to the Tarski's
World exercises used in logic courses~\cite{barwise}.}
\label{fig:tarski}
\end{figure}

\begin{exercise}[easy]\label{exr:negate-quant}
Negate ``every contestant solved some problem'' from Example~\ref{ex:quant-swap},
first in English, then in symbols. Check against world 3 of the listing.
\end{exercise}

\section{Negating statements mechanically}\label{sec:negation}
\index{negation!of quantified statements}

Negation is the highest-frequency operation in proof work: every contradiction,
every contrapositive, every disproof starts by negating something. Do it by
mechanical rules, not by feel.

\begin{principle}[Negate from the outside in]
To negate a statement, push $\lnot$ inward past each operator, flipping as you go:
\[
\lnot\forall \to \exists\lnot,\qquad
\lnot\exists \to \forall\lnot,\qquad
\lnot(P\land Q) \to \lnot P \lor \lnot Q,
\]
\[
\lnot(P\lor Q) \to \lnot P \land \lnot Q,\qquad
\lnot(P \Rightarrow Q) \to P \land \lnot Q.
\]
When $\lnot$ reaches an atomic predicate it stops and negates it directly.
\end{principle}

The last rule follows from $P\Rightarrow Q \equiv \lnot P \lor Q$: negating gives
$P \land \lnot Q$. In English: the only way ``if $P$ then $Q$'' fails is a case where
the hypothesis happens and the conclusion does not.

\begin{example}\label{ex:negate-chain}
Negate $\forall x\,\exists y\, R(x,y)$.
\par\smallskip
\textbf{Strategy.} Work left to right. Each quantifier flips once, the innermost
predicate negates.
\par\smallskip
\begin{solution}
$\lnot\bigl(\forall x\,\exists y\,R(x,y)\bigr)
\equiv \exists x\,\lnot\bigl(\exists y\,R(x,y)\bigr)
\equiv \exists x\,\forall y\,\lnot R(x,y)$.
One witness $x$ for which \emph{every} $y$ fails is what refutes the original.
\end{solution}
\end{example}

\begin{tripwire}
A very common slip negates $\forall x\,\exists y$ as $\exists x\,\exists y$: both
quantifiers flipped, but the order was left untouched. The correct negation is
$\exists x\,\forall y$. The $\exists$-witness comes \emph{first} in the negation,
and it must defeat \emph{all} $y$.
\end{tripwire}

\begin{figure}[ht]\centering
% Figure: pushing a negation through a quantifier chain.
% Source for Figure~\ref{fig:negpush}. TikZ only; regenerable by compiling main.tex.
\begin{tikzpicture}[
  box/.style={draw,inner sep=4pt,font=\small,rounded corners=1pt},
  lbl/.style={font=\footnotesize\sffamily,text=slate,align=left},
  arr/.style={->,thick,accent}
]
\node[box] (s0) at (0,0) {$\lnot\,\forall x\,\exists y\,P(x,y)$};
\node[box] (s1) at (0,-1.15) {$\exists x\,\lnot\,\exists y\,P(x,y)$};
\node[box] (s2) at (0,-2.3) {$\exists x\,\forall y\,\lnot P(x,y)$};
\draw[arr] (s0) -- (s1);
\draw[arr] (s1) -- (s2);
\node[lbl] at (4.6,0) {outermost symbol first:};
\node[lbl] at (4.6,-1.15) {$\lnot\forall$ becomes $\exists$,};
\node[lbl] at (4.6,-2.3) {$\lnot\exists$ becomes $\forall$, negation lands on $P$.};
\node[lbl,anchor=west] at (-3.4,-3.15) {Each quantifier flips exactly once; order never changes.};
\end{tikzpicture}

\caption{Negation pushed inward through a two-quantifier chain. Notice:
each quantifier flips exactly once, the order is untouched, and the
negation lands on the innermost predicate.}
\label{fig:negpush}
\end{figure}

% from-zero:naive-neg:begin
\begin{example}[First attempt: negating by flipping every word]\label{ex:naive-negation}
A teacher says: ``Every student passed the exam.'' You want to say this is false. The
natural first attempt is ``Every student failed the exam.''
\par\smallskip
\textbf{What goes wrong.} Suppose $20$ students sat the exam and $19$ passed. The teacher's
sentence is false (one student did not pass), yet ``every student failed'' is also false
($19$ passed). The attempted negation is false \emph{in a case where the original is false too},
so it cannot be the negation. In symbols, with $P(s)$ meaning ``$s$ passed'', the original is
$\forall s\,P(s)$. The correct negation is $\exists s\,\neg P(s)$: ``there is a student who did
not pass.'' Flipping the quantifier \emph{and} the statement inside it is the rule of
Section~\ref{sec:negation}; flipping only the inside is the commonest error in this book.
\end{example}
% from-zero:naive-neg:end

\section{English and symbols}\label{sec:english}
\index{translation!English--symbols}

Contest statements arrive in English; proofs need symbols. The translation step is
where most wrong readings happen, so we make it explicit. The usual hiding places:
``a'', ``an'' and ``the'' often mean $\exists$; ``every'', ``each'', ``any'' and
``no matter which'' mean $\forall$; ``some'' and ``at least one'' mean $\exists$;
``no'' means $\lnot\exists$ (equivalently $\forall\dots\lnot$); ``only'' reverses
arrows --- ``$P$ only if $Q$'' is $P \Rightarrow Q$, not $Q \Rightarrow P$.

\begin{figure}[ht]\centering
% Figure: English-to-symbols decision flow for quantified statements.
% Source for Figure~\ref{fig:flow}. TikZ only; regenerable by compiling main.tex.
\begin{tikzpicture}[
  q/.style={draw,inner sep=4pt,font=\footnotesize,rounded corners=1pt,align=center},
  a/.style={draw,inner sep=4pt,font=\footnotesize,align=center,fill=rulegray!8},
  arr/.style={->,thick,rulegray},
  lab/.style={font=\tiny\sffamily,text=slate}
]
\node[q] (start) at (0,0) {``\emph{Every / All / Each}''?};
\node[q] (some) at (0,-1.5) {``\emph{Some / There is / At least one}''?};
\node[q] (onlyif) at (0,-3.0) {``\emph{Only if / unless / whenever}''?};

\node[a] (fa) at (6.2,0) {$\forall x\,(\mathrm{scope}(x)\Rightarrow \mathrm{claim}(x))$\\
\tiny restrict the domain, then require};
\node[a] (ex) at (6.2,-1.5) {$\exists x\,(\mathrm{scope}(x)\land \mathrm{claim}(x))$\\
\tiny restrict the domain, then exhibit};
\node[a] (oi) at (6.2,-3.0) {reverse the arrow:\\
\tiny ``$P$ only if $Q$'' $=$ $P\Rightarrow Q$};

\draw[arr] (start) -- (fa);
\draw[arr] (some) -- (ex);
\draw[arr] (onlyif) -- (oi);

\node[lab,anchor=west] at (-4.9,-4.15) {Universal quantifiers pair with $\Rightarrow$; existential with $\land$.
Swapping the connective is the single most common mistranslation.};
\end{tikzpicture}

\caption{The three recurring English patterns and their connectives.
Notice the left column: universals take $\Rightarrow$ because you restrict
the domain and then \emph{require} the claim; existentials take $\land$
because you restrict the domain and then \emph{exhibit} a member.}
\label{fig:flow}
\end{figure}

\begin{example}[Goldbach's claim]\label{ex:goldbach}
Write in symbols, then negate: \emph{every even integer greater than $2$ is a sum of
two primes}.
\par\smallskip
\textbf{Strategy.} Outer structure first. The sentence is a ``for all'' over even
integers with a condition ($>2$, even), so it is $\forall n \bigl((n \text{ even}
\land n>2) \Rightarrow \dots\bigr)$ --- a universal over a restricted domain, which
always means a $\forall$ wrapped around an implication. The inner claim ``is a sum
of two primes'' is existential: $\exists p\,\exists q$ with $p,q$ prime and
$p+q=n$. Then negate from the outside in.
\par\smallskip
\begin{solution}
Let $E(n)$, $P(n)$ abbreviate ``$n$ is even'' and ``$n$ is prime''. The claim is
\[
\forall n\in\Z\;\Bigl( \bigl(E(n)\land n>2\bigr) \Rightarrow
\exists p\,\exists q\;\bigl(P(p)\land P(q)\land p+q=n\bigr)\Bigr).
\]
Negating: $\exists n\in\Z$ such that $E(n)\land n>2$ and
$\forall p\,\forall q\,\lnot\bigl(P(p)\land P(q)\land p+q=n\bigr)$ ---
\emph{some even $n>2$ that is not a sum of two primes}. One counterexample kills the
claim; no amount of examples proves it. This is \emph{Goldbach's
conjecture}\index{Goldbach conjecture}: stated 1742, still unproved as of
2026 --- yet mechanically, its negation took two lines.
\end{solution}
\end{example}

\begin{wins}[A translation is done when\dots]
every English word has become a named object, a quantifier, or a connective ---
and the direction of every ``only if'' and every ``unless'' has been checked on a
case where it matters. If a quantifier's universe is left implicit in English, name
it in symbols.
\end{wins}

\begin{remark}
Competition problems are dense quantified statements wearing prose clothes.
AMC~12~2024~A~P22 is casework counting organised around a closed-loop (universal)
constraint; AMC~12~2024~B~P5 is a sign-flip problem whose engine is magnitude
bookkeeping --- each flip changes the total by twice the flipped term, and the
parity of the sum never moves. Both are worth reading on the archive~
\cite{aops-amc} once you can spot the quantified claims inside them.
\end{remark}

\begin{example}[``There is a unique\ldots'']\label{ex:unique}
Translate: \emph{there is a unique prime $p$ such that $p{+}2$ is prime}. Then
evaluate the statement.
\par\smallskip
\textbf{Strategy.} ``Unique'' is two claims AND-ed: existence ($\exists p$) and
at-most-one (any two candidates coincide). Write both halves, then hunt for the
witness --- every prime except $2$ is odd, and an odd $p$ makes $p{+}2$ odd only
in the right way, so check small candidates directly.
\par\smallskip
\begin{solution}
With $P(n)$ = ``$n$ is prime'', uniqueness is
\[
\exists p\,\Bigl(P(p)\land P(p{+}2)\Bigr)\;\land\;
\forall q\,\forall r\,\bigl(\bigl(P(q)\land P(q{+}2)\land P(r)\land P(r{+}2)\bigr)
\Rightarrow q=r\bigr).
\]
Existence: $p=3$ works ($3,5$ both prime). At-most-one fails immediately:
$q=3$ and $r=5$ both satisfy the property but $3\neq 5$ --- $(5,7)$ is another
prime pair. The English statement is \emph{false}: there are (at least) two such
primes, so the ``unique'' claim dies on the second half, not the first. The
English ``unique'' is exactly the kind of word that silently adds a second
quantified clause --- always expand it.
\end{solution}
\end{example}

\begin{exercise}[medium]\label{exr:translate}
Write in symbols over $\Z$: ``exactly one of the numbers $a,b,c$ is even''. Then
decide: is the resulting statement, quantified over all $a,b,c\in\Z$, true, false,
or an open formula? Justify.
\end{exercise}

\section{Parity and primes: a first look}\label{sec:parity}
\index{parity}\index{even}\index{odd}\index{prime}\index{divides}

The grammar of this chapter is now complete; here is some content to practise it on.
Two definitions run through the whole book.

\begin{definition}[Even, odd, divides]
An integer $n$ is \emph{even} when $\exists k\in\Z\;(n=2k)$, and \emph{odd} when
$\exists k\in\Z\;(n=2k+1)$. For integers $d,m$ we say \emph{$d$ divides $m$} when
$\exists k\in\Z\;(m=d\cdot k)$.
\end{definition}

\begin{definition}[Prime]\index{composite}
An integer $n>1$ is \emph{prime} when its only positive divisors are $1$ and $n$:
in symbols, $\forall d\in\Z^{+}\;\bigl(d \text{ divides } n \Rightarrow
(d=1 \lor d=n)\bigr)$. An $n>1$ that is not prime is \emph{composite}.
\end{definition}

Both definitions are quantified statements. ``$n$ is prime'' is a
$\forall\Rightarrow$ statement --- the expensive kind, which is why proving primality
means excluding every candidate divisor, while disproving it needs one witness.
Already you can classify the work: ``every prime $>2$ is odd'' is
$\forall p\;\bigl((P(p)\land p>2)\Rightarrow \text{odd}(p)\bigr)$; a proof (Chapter~2)
must take an arbitrary such $p$ and derive oddness, while one even prime $>2$ would
kill it. The claims and their forms:
\begin{center}
\begin{tabular}{lcc}
\toprule
Claim & Shape & To refute it you need\\
\midrule
every prime $>2$ is odd & $\forall\Rightarrow$ & one even prime $>2$\\
every even $n>2$ is composite & $\forall\Rightarrow$ & one even prime $>2$\\
there are infinitely many primes & $\forall n\,\exists p\;(p>n \land P(p))$ & a largest prime\\
some even $n$ is a sum of two primes & $\exists$ & check all --- impossible\\
\bottomrule
\end{tabular}
\end{center}
Notice the last row: a bare $\exists$ claim can almost never be \emph{refuted} by
search, only verified. That asymmetry is why existence claims and universal claims
play different roles in problems.

Parity is governed by a small arithmetic. You may verify all four rows yourself ---
each is a two-line calculation once ``even'' is read as $2k$:
\begin{center}
\begin{tabular}{cc|cc}
\toprule
$x$ & $y$ & $x+y$ & $xy$\\
\midrule
even & even & even & even\\
even & odd & odd & even\\
odd & even & odd & even\\
odd & odd & even & odd\\
\bottomrule
\end{tabular}
\end{center}
This table is the entire engine behind parity arguments --- and it is also worth
knowing where it does \emph{not} apply. AMC~12~2024~B~P5 asks for the fewest
$+$-to-$-$ sign flips in $1+3+\cdots+99$ that make the sum negative; flipping a
term $t$ changes the total by $2t$, so the sum's parity is an invariant and the
real engine is magnitude --- how few of the largest terms sum past half of
$2500$. Proofs that this table is correct come in Chapter~2; the shape to
notice is \emph{$\forall$ pairs of integers, a finite check of four cases} ---
a universal statement with a finite universe in disguise.

\begin{exercise}[medium]\label{exr:findn}
Find all positive integers $n$ for which the statement ``$n$ divides $n+6$'' is
true. (Use the division definition above: $d$ divides $m$ when $m=d\cdot k$ for
some integer $k$.)
\end{exercise}

\section{Chapter summary}
\begin{table}[ht]\centering
\begin{tabularx}{\linewidth}{@{}lcYY@{}}
\toprule
\tablehead\thead{Symbol} & \thead{Name} & \thead{English forms} & \thead{Negation}\\
\midrule
$P\land Q$ & conjunction & and; both; but & $\lnot P\lor\lnot Q$\\
$P\lor Q$ & disjunction & or (inclusive); at least one of & $\lnot P\land\lnot Q$\\
$\lnot P$ & negation & not; it is false that & $P$\\
$P\Rightarrow Q$ & implication & if $P$ then $Q$; $P$ only if $Q$; $Q$ whenever $P$ & $P\land\lnot Q$\\
$P\iff Q$ & biconditional & iff; exactly when; equivalent & $(P\land\lnot Q)\lor(\lnot P\land Q)$\\
$\forall x\,P(x)$ & universal & every; each; any; no matter which & $\exists x\,\lnot P(x)$\\
$\exists x\,P(x)$ & existential & some; there is; at least one & $\forall x\,\lnot P(x)$\\
$\exists! x\,P(x)$ & unique existence & exactly one & $\forall x\,\lnot P(x)\ \lor$ two distinct witnesses\\
\bottomrule
\end{tabularx}
\caption{The grammar of the book on one page. The last column is the mechanical
negation rule --- the most-used column.}
\label{tab:grammar}
\end{table}

\begin{table}[ht]\centering
\begin{tabularx}{\linewidth}{@{}lYY@{}}
\toprule
\tablehead\thead{Claim shape} & \thead{To establish it} & \thead{To refute it}\\
\midrule
$\forall x\,P(x)$ & handle every $x$ (proof, Chapter~2) & one counterexample\\
$\exists x\,P(x)$ & one witness & rule out every $x$\\
$P\Rightarrow Q$ & assume $P$, derive $Q$; or assume $\lnot Q$, derive $\lnot P$ & one case with $P$ true, $Q$ false\\
$P\iff Q$ & both directions, separately & one direction fails\\
$\forall x\,\exists y$ & give each $x$ its own $y$ & one $x$ defeating all $y$\\
$\exists y\,\forall x$ & one $y$ serving all $x$ & for each $y$, some $x$ it fails\\
\bottomrule
\end{tabularx}
\caption{What a proof or a refutation owes you, per claim shape.}
\label{tab:jobs}
\end{table}

\section{Common mistakes}
\index{common mistakes}
\begin{itemize}
\item \textbf{Asserting the converse.} From ``if $n$ is divisible by $6$ then $n$ is
even'' and ``$n$ is even'' you get nothing --- $n=10$ satisfies the conclusion and
not the hypothesis. The converse $Q\Rightarrow P$ is a \emph{different} claim
(Figure~\ref{fig:family}); proving it is a separate job.
\item \textbf{Negating $\forall\exists$ as $\exists\exists$.} The negation of
``every row has a mark'' is ``\emph{some} row has \emph{no} marks''
($\exists x\,\forall y$), not ``some row has some mark''. Flip each quantifier once
and keep the order.
\item \textbf{Reading ``or'' as exclusive.} ``$x$ is even or $x$ is prime'' is true
of $x=2$ twice over; disjunction does not promise exactly one branch. When a
statement means exactly one, it says so --- translate that as two claims
($\lor$ plus $\lnot\land$), as in Exercise~\ref{exr:translate}.
\item \textbf{Negating an implication as another implication.}
$\lnot(P\Rightarrow Q)$ is $P\land\lnot Q$, not $P\Rightarrow\lnot Q$. A
counterexample to ``if $n$ is prime then $n$ is odd'' is a prime that is
\emph{not} odd ($n=2$) --- an actual object, not a reversed claim.
\item \textbf{Marking vacuously true claims false.} ``Every even prime greater
than $2$ is divisible by $5$'' is \emph{true}: there are no even primes above
$2$, so the implication never gets a test case. Vacuous truth feels like a
cheat; it is the correct reading.
\item \textbf{Assuming one witness serves two existential claims.}
From $\exists x\,P(x)$ and $\exists x\,Q(x)$ you cannot conclude
$\exists x\,\bigl(P(x)\land Q(x)\bigr)$ --- the $x$'s may be different objects
(some number is even; some number is odd; no number is both). Existentials
only merge when one statement produces the witness the other needs.
\end{itemize}

\section{Additional exercises}
\begin{exercise}[easy]
Write the truth table of $(P \lor Q) \land \lnot(P \land Q)$ and compare it with
``exactly one of $P,Q$''. This is the \emph{exclusive or}.
\end{exercise}
\begin{exercise}[easy]
Negate mechanically: (a) $\exists x\,\forall y\,R(x,y)$;
(b) $\forall x\,\bigl(P(x)\Rightarrow Q(x)\bigr)$;
(c) ``there is no largest prime'' after writing it in symbols.
\end{exercise}
\begin{exercise}[easy]
Over $\Z$, is $\exists x\,\forall y\;(x\cdot y = y)$ true or false? Give a witness
or a refutation, and state the claim in English.
\end{exercise}
\begin{exercise}[easy]
Classify each as statement, open formula, or neither: (a) $n^2$ is even;
(b) $n$ is even for every $n$; (c) every multiple of $4$ is even; (d) prove that
$n$ is even.
\end{exercise}
\begin{exercise}[medium]
Write in symbols and determine truth over $\Z$: (a) ``for every $x$ there is a $y$
with $x+y=0$''; (b) ``there is a $y$ such that for every $x$, $x+y=0$''.
Explain the difference in one sentence.
\end{exercise}
\begin{exercise}[medium]
Prove or disprove: if $x^2$ is divisible by $4$, then $x$ is divisible by $4$.
(If disproving, exhibit the smallest counterexample.)
\end{exercise}
\begin{exercise}[medium]
Negate the $\varepsilon$--$\delta$ statement you will meet in analysis:
$\forall \varepsilon>0\,\exists\delta>0\,\forall x\,
(|x-a|<\delta \Rightarrow |f(x)-f(a)|<\varepsilon)$.
Your answer should begin ``$\exists \varepsilon>0$ such that\dots''.
\end{exercise}
\begin{exercise}[medium]
On the set $\{1,2,3\}$, define a relation $R(x,y)$ such that
$\forall x\,\exists y\,R(x,y)$ is true and $\exists y\,\forall x\,R(x,y)$ is false;
then one where both are true. Draw both as grids of marks like
Figure~\ref{fig:quant-order}.
\end{exercise}
\begin{exercise}[hard]
Define $\exists!x\,P(x)$ two ways: (i) $\exists x\,\bigl(P(x)\land \forall y\,
(P(y)\Rightarrow y=x)\bigr)$ and (ii) $\exists x\,P(x)\;\land\;\forall x\,\forall y\,
\bigl(P(x)\land P(y) \Rightarrow x=y\bigr)$. Prove the two readings agree: each
implies the other. (This is the shape of every ``unique solution'' claim in the
book.)
\end{exercise}
\begin{exercise}[hard]
Extend \texttt{code/verify\_quantifiers.py} into a tiny verifier: input a finite
grid of marks and a claim of the form $\forall x\,\exists y$, $\exists y\,\forall x$,
or their negations, and have it print the verdict. Test it on the three worlds of
the listing and on Figure~\ref{fig:quant-order}'s two grids.
\end{exercise}
\begin{exercise}[easy]
Sort these into equivalent pairs: $P\Rightarrow Q$; $\lnot Q\Rightarrow\lnot P$;
$Q\Rightarrow P$; $\lnot P\Rightarrow\lnot Q$; $\lnot P\lor Q$; $\lnot(P\land\lnot Q)$.
Which pairs are provably equal by truth table?
\end{exercise}
\begin{exercise}[medium]
Decide the truth value of each claim over $\Z$, and say whether any is vacuously
true: (i) every negative perfect square is divisible by $4$; (ii) every integer
$n>1$ has a divisor $d$ with $1<d\le\sqrt{n}$ or $n$ is prime; (iii) some
composite $n$ has no prime factor at most $\sqrt{n}$.
\end{exercise}
\begin{exercise}[medium]
Write ``there is a unique even prime'' in symbols using $\exists!x\,P(x)$
expanded as in the earlier hard exercise, then prove the statement true by
exhibiting the witness and showing uniqueness.
\end{exercise}
\begin{exercise}[medium]
A classmate translated ``no even number above $2$ is prime'' as
$\forall n\,\bigl(\text{even}(n)\land n>2 \Rightarrow \lnot\text{prime}(n)\bigr)$
--- correct --- but then negated it as $\exists n\,(\text{even}(n)\land n>2
\Rightarrow \text{prime}(n))$. Find the two errors and write the real negation
in English.
\end{exercise}
\begin{exercise}[hard]
Prove by truth table that $\lnot(P\Rightarrow Q)\equiv P\land\lnot Q$, then use it
to negate, in one line each: (i) ``if $n$ is even then $n^2$ is divisible by $4$'';
(ii) ``whenever the algorithm halts, its output is sorted''. Give the English
reading of each negation.
\end{exercise}
\begin{exercise}[hard]
``No'' hides a negated existential. Write ``no prime above $10$ ends in $5$''
three ways: as $\lnot\exists$, as $\forall\ldots\lnot$, and as a single
implication over $\Z$. Verify all three agree on a truth-table-style check
over $n\in\{2,\dots,20\}$ (by hand or code).
\end{exercise}

\section*{Further reading}
\begin{itemize}
\item R.~Hammack, \emph{Book of Proof}, 3rd ed., Ch.~2 ---
\url{https://richardhammack.github.io/BookOfProof/} --- free PDF; the
friendliest full treatment of connectives and quantifiers, with a large
exercise bank pitched exactly at this chapter.
\item MIT OpenCourseWare, \emph{6.042J Mathematics for Computer Science}
(Fall 2010), Ch.~1 --- \url{https://ocw.mit.edu/6-042JF10} ---
the course this book's proof-first order follows; Ch.~1 covers
propositions, implication and proof patterns.
\item E.~Lehman, F.~T.~Leighton, A.~R.~Meyer, \emph{Mathematics for Computer
Science} (the 6.042J text), \S1 ---
\url{https://courses.csail.mit.edu/6.042/spring18/mcs.pdf} --- free PDF;
denser than Hammack, worth it for the ``logic in programs'' angle.
\item Stanford CS103, \emph{Mathematical Foundations of Computing} course
reader ---
\url{https://web.stanford.edu/class/cs103/} --- free reader + problem sets;
its logic units are the best drills on translation and negation.
\item The Open Logic Project, \emph{forall x: Calgary} ---
\url{https://openlogicproject.org/} --- free intro-logic text; go here when
you want the formal (Fitch-style) counterpart of what this chapter teaches
informally.
\item CMU 15-151, \emph{Mathematical Foundations for Computer Science}
(Great Ideas in Theoretical CS) ---
\url{https://www.cs.cmu.edu/~15281/} --- lecture notes by Anil Ada et al.;
the early chapters put propositional logic to work on real problems fast.
\end{itemize}

\begin{rulebox}[Watch instead]
\begin{itemize}
\item \textbf{MIT OpenCourseWare --- ``Lec 1 | MIT 6.042J''}, Tom Leighton,
44:09, \url{https://www.youtube.com/watch?v=L3LMbpZIKhQ} --- covers
\S1.1--1.3: propositions, truth tables, implication, and the first proofs.
(Verified 2026-10-11.)
\item \textbf{Neso Academy --- ``Logical Operators: Negation, Conjunction \&
Disjunction''}, 6:09,
\url{https://www.youtube.com/watch?v=6kYngPvoGxU} --- covers \S1.1--1.2: the
connectives and their truth tables, the quickest possible refresher.
(Verified 2026-10-11.)
\item \textbf{Mathematical Visual Proofs --- ``Predicates and Quantifiers
[Discrete Math Class]''}, $\sim$10:00,
\url{https://www.youtube.com/watch?v=0rvKhma-3f4} --- covers \S1.4--1.5:
predicates, quantifier order, and the negation laws that match
Figure~\ref{fig:negpush}. (Verified 2026-10-11.)
\end{itemize}
\end{rulebox}

\chapternotes
The order --- connectives, truth tables, implication family, then quantifiers and
mechanical negation --- follows Velleman's first two chapters, the clearest
prose treatment of this material~\cite{velleman}, and Hammack's logic chapter, whose
exercise mix (translate, negate, truth-table, find-the-error) shaped ours~\cite{hammack}.
The proof-first ordering across the whole book follows MIT 6.042J~\cite{mcs}.
The finite-world evaluation of Example~\ref{ex:tarski} is in the tradition of
Barwise and Etchemendy's \emph{Tarski's World}~\cite{barwise}; Cut-the-Knot's logic
and fallacy pages~\cite{cuttheknot} supply many of the standard traps catalogued in
Common mistakes. Problem citations (AMC~12~2024~A~P22, AMC~12~2024~B~P5) point to
the AoPS archive~\cite{aops-amc}; read the statements there --- they are referenced,
never reproduced.
